\scanpage{75}
3) \uline{Try to find reasonable definitions in characteristic}2, \uline{or preferably independent of the characteristic}. A really satisfactory theory should work over every prescheme $S$, and conversely. I am far from having clear ideas on the question, and confine myself to vague remarks:

a) When one is over a field of \inserted{arbitrary} characteristic \deleted{2}, one is probably obliged to take degenerate quadratic forms into account, and to define a \deleted{\unsure{r}}ring $Q'(k)$ as defined by the classes of quadratic forms, degenerate or not. \handwritten{[}In the case of a field of \deleted{\illegible{1}} characteristic $\handwritten{\ne}2$, this would be $Q(k)+\underline{Z}$, where the \deleted{second factor $\underline{Z}$ is a} multiplication is given by $(\unsure{g|f},m).(g,n)=(fg,\deg(f)n+\deg(g)m+mn)$, the extra factor $\underline{Z}$ corresponding to the zero forms, on vector spaces of arbitrary rank\handwritten{]}.

b) In arbitrary characteristic, keep in mind the distinction between the group scheme $(\underline{Z}/2\underline{Z})_S$ and the group scheme $(\handwritten{\mu_2})_S$ over $S$, the latter being the kernel of squaring in the multiplicative group $\underline{G}_{m\,S}$. If $\unsure{S|A}$ is local with ring $A$, then the exact cohomology sequence gives $H^{\handwritten{1}}(S,\handwritten{\mu_2})=A^{\handwritten{*}}/A^{\handwritten{*2}}$, $H^{\handwritten{2}}(S,\handwritten{\mu_2})={}_{\handwritten{2}}H^{\handwritten{2}}(A,\underline{G}_m)$, elements annihilated by 2 in the Brauer group, on the other hand $H^{\handwritten{1}}(S,\underline{Z}/2\underline{Z})=A/\handwritten{\wp}A$ when $2A=0$. One must therefore expect to find Stiefel-Whitney classes of mixed type, involving the two neighbouring coefficient groups at once. This remark becomes more precise when one looks at the orthogonal group scheme\handwritten{s}, over $\underline{Z}$ say: in even dimension, it is the extensions of $\underline{Z}/2\underline{Z}$ by a group scheme with connected fibers (the special orthogonal scheme), for example in dimension 2 it is \deleted{the extension} a semidirect product of $\underline{Z}/2\underline{Z}$ by $\underline{G}_m$ on which it acts by symmetry; in odd dimension, it is an extension of $\handwritten{\mu_2}$ by the special orthogonal group, for example in dimension 1, it is $\handwritten{\mu_2}$. Hence depending on the parity, as first ``discrete'' invariant of a principal orthogonal bundle, a 1-class with group $\underline{Z}/2\underline{Z}$,or $\handwritten{\mu_2}$.
% archive sheet number: 74
